% Generated draft: requires mathematical and source review.
\documentclass[11pt]{article}
\usepackage[margin=25mm]{geometry}
\usepackage{fontspec}
\setmainfont{Latin Modern Roman}
\usepackage{amsmath,amssymb,mathtools,unicode-math}
\setmathfont{Latin Modern Math}
\usepackage{enumitem,needspace,xurl,hyperref}
\setlength{\parindent}{0pt}
\setlength{\parskip}{5pt}
\setlength{\emergencystretch}{3em}
\newcommand{\N}{\mathbb{N}}
\newcommand{\Z}{\mathbb{Z}}
\newcommand{\Q}{\mathbb{Q}}
\newcommand{\R}{\mathbb{R}}
\newcommand{\C}{\mathbb{C}}
\newcommand{\F}{\mathbb{F}}
\newcommand{\A}{\mathbb{A}}
\newcommand{\PP}{\mathbb{P}}
\newcommand{\E}{\mathbb{E}}
\newcommand{\eps}{\varepsilon}
\newcommand{\dd}{\,\mathrm{d}}
\begin{document}
\section*{UnsolvedMath ID: 20002905}
\textbf{Source problem code:} catalogue entry
\subsection{Definitions and mathematical statement}
The workshop asks how $k$-Schur functions relate to $k$-level-restricted Schur functions, specifically whether all level-restricted structure constants occur among the $k$-Littlewood--Richardson coefficients. The source suggests the Goodman--Wenzl root-of-unity Hecke-algebra work as relevant. It does not define its rank, level or quotient conventions.

The supplied aid makes a type-$A$ WZW interpretation. Let $\ell\geq2$, let the Lie rank be $r=\ell-1$, and let $L\geq1$ be the level. The relevant Schur cutoff is
\[
K=r+L=\ell+L-1,
\]
rather than simply $L$. Label the $\operatorname{SU}(\ell)$ level-$L$ fusion basis by partitions $\lambda$ with at most $\ell-1$ rows and $\lambda_1\leq L$. If $[\lambda][\mu]=\sum_\nu N_{\lambda,\mu}^{\nu}[\nu]$, the $N$ are its fusion coefficients. At parameter $1$, define the homogeneous $K$-Schur coefficients by
\[
s_\alpha^{(K)}s_\beta^{(K)}=\sum_\gamma c_{\alpha,\beta}^{\gamma,K}s_\gamma^{(K)}.
\]
For a partition $\lambda$, write $\lambda'$ for its transpose.

The aid attributes to Lapointe--Morse a root-of-unity Hecke quotient in which a $K$-Schur function maps either to its ordinary Schur class or to zero, yielding the desired realization of level-restricted constants. In the transposed fusion convention it reports
\[
N_{\lambda,\mu}^{\nu}=c_{\lambda',\mu'}^{(\ell^d,\nu'),K},\qquad
d=\frac{|\lambda|+|\mu|-|\nu|}{\ell},
\]
when $d$ is a nonnegative integer. Here $(\ell^d,\nu')$ means prepending $d$ full rows of length $\ell$ to $\nu'$. The charge and admissibility conventions for zero coefficients when no such lift exists must be matched to the precise quotient theorem; they are not fully stated in this record.

The aid's homogeneity observation is that the fusion relation $s_{(\ell)}=1$ erases these full rows, while the original homogeneous product forces
\[
|\lambda'|+|\mu'|=d\ell+|\nu'|.
\]
Thus there is at most one nonnegative full-row lift of the indicated type that can contribute. It is the lifted index, rather than necessarily the unchanged index $\nu'$, that recovers the fusion coefficient. For its example $\operatorname{SU}(2)$ at level $2$, the cutoff is $K=3$; the coefficient of the unit in the square of the fundamental class is $1$, corresponding to $c_{(1),(1)}^{(2),3}=1$ with $d=1$, whereas the homogeneous coefficient with output the empty partition is zero.

The relationship is therefore a quotient and index-lifting relationship in this interpretation. The infinite graded $K$-Schur algebra cannot literally be the finite fusion ring with the same grading, since the latter identifies a positive-degree full-row class with the unit. Exact quotient ideals, parameter specialization and the Hecke-to-fusion dictionary remain necessary review data.
\subsection{Short English statement}
The supplied type-$A$ dictionary realizes level-restricted fusion constants as selected $K$-Schur coefficients, with $K$ equal to rank plus level. Homogeneity determines a necessary full-row lift of the output partition, so the relation is a quotient rather than literal equality of the two algebras.
\subsection{Sources}
\begin{itemize}
\item Local UnsolvedMath registry, permanent ID $20002905$, source code \texttt{catalogue entry}; canonical \texttt{aim-special-functions-notes.json}, \texttt{notes[9]}.
\item AIM workshop \emph{Generalized Kostka polynomials}, Problem (10): \url{https://aimath.org/WWN/kostka/prob2.pdf}. No original remarks or separate literature field are supplied.
\item Goodman and Wenzl, paper on Iwahori--Hecke algebras of type $A$ at roots of unity, \emph{J. Algebra} \textbf{215} (1999), 694--734, as given in the source. The problem says this paper may be relevant.
\item The supplied aid attributes the $K$-Schur quotient theorem to Lapointe--Morse without a full bibliographic reference or quotient presentation.
\item Research attempt 1: \texttt{partial\_result}, status \texttt{solved}, classification \texttt{SOLVED-IN-LITERATURE}, rigor \texttt{proved}, novelty confidence low. Its forced-lift proposition and literature interpretation are machine-generated aid, not peer reviewed, and require independent verification before citation.
\end{itemize}
\end{document}
